% ═══════════════════════════════════════════════════════════════
% AB-07c: Querer y Poder - Modalverben beim Einkaufen
% Abgeleitet aus LE-07c (Score 9.3/10)
% ═══════════════════════════════════════════════════════════════

\documentclass[11pt, a4paper]{article}

% ═══════════════════════════════════════════════════════════════
% SW-MODUS TOGGLE
% ═══════════════════════════════════════════════════════════════
% Setze \swmodustrue für Schwarz-Weiß-Druck
\newif\ifswmodus
\swmodusfalse    % Standard: Farbe

\usepackage[utf8]{inputenc}
\usepackage[T1]{fontenc}
\usepackage[ngerman]{babel}
\usepackage{helvet}
\renewcommand{\familydefault}{\sfdefault}

\usepackage[margin=2cm]{geometry}
\usepackage{enumitem}
\usepackage{tabularx}
\usepackage{booktabs}

\usepackage{xcolor}
\usepackage{tcolorbox}
\tcbuselibrary{skins}

\usepackage{fancyhdr}
\usepackage{lastpage}
\usepackage{needspace}

% Farben - Basis
\definecolor{spanisch-color}{HTML}{c41e3a}
\definecolor{grau-color}{HTML}{888888}
\definecolor{hellgrau-color}{HTML}{f5f5f5}
\definecolor{orange-color}{HTML}{b35c00}

% SW-Farben (druckerfreundlich)
\definecolor{sw-dark}{HTML}{000000}
\definecolor{sw-gray}{HTML}{666666}
\definecolor{sw-white}{HTML}{ffffff}

% Bedingte Farbzuweisung
\ifswmodus
    \colorlet{spanisch}{sw-dark}
    \colorlet{grau}{sw-gray}
    \colorlet{hellgrau}{sw-white}
    \colorlet{orange}{sw-dark}
\else
    \colorlet{spanisch}{spanisch-color}
    \colorlet{grau}{grau-color}
    \colorlet{hellgrau}{hellgrau-color}
    \colorlet{orange}{orange-color}
\fi

\newcommand{\fachfarbe}{spanisch}

% Variablen
\newcommand{\thema}{Querer y Poder -- Modalverben beim Einkaufen}
\newcommand{\sequenz}{07}
\newcommand{\block}{c}
\newcommand{\fach}{Spanisch}
\newcommand{\klasse}{7}
\newcommand{\maxpunkte}{24}

% Kopf-/Fußzeile
\pagestyle{fancy}
\fancyhf{}
\renewcommand{\headrulewidth}{0.4pt}
\renewcommand{\footrulewidth}{0.4pt}
\lhead{\fach{} -- Klasse \klasse}
\chead{AB-\sequenz\block}
\rhead{}
\lfoot{}
\rfoot{Seite \thepage\ von \pageref{LastPage}}

% Befehle
\newcommand{\punktebox}[1]{\hfill\fbox{\small \_/#1}}
\newcommand{\luecke}[1]{\rule{#1}{0.4pt}}
\newcommand{\es}[1]{\textcolor{\fachfarbe}{\textbf{#1}}}

% Merkkasten (SW-kompatibel)
\newtcolorbox{merkkasten}[1][]{
    colback=hellgrau,
    colframe=\fachfarbe,
    colbacktitle=white,
    coltitle=\fachfarbe,
    boxrule=1pt,
    arc=0pt,
    left=8pt, right=8pt, top=8pt, bottom=8pt,
    fonttitle=\bfseries,
    attach boxed title to top left={yshift=-2mm, xshift=5mm},
    boxed title style={boxrule=0pt, colback=white},
    title=#1
}

% Hinweis (SW-kompatibel)
\newtcolorbox{hinweis}{
    colback=white,
    colframe=orange,
    colbacktitle=white,
    coltitle=orange,
    boxrule=1pt,
    arc=0pt,
    left=5pt, right=5pt, top=5pt, bottom=5pt,
    fonttitle=\bfseries,
    attach boxed title to top left={yshift=-2mm, xshift=5mm},
    boxed title style={boxrule=0pt, colback=white},
    title=Hinweis
}

% Aufgabe
\newenvironment{aufgabe}[1]{%
    \needspace{5\baselineskip}%
    \nopagebreak[4]%
    \vspace{10pt}
    \noindent\textbf{\textcolor{\fachfarbe}{Aufgabe #1}}
    \nopagebreak[4]%
    \vspace{5pt}
    \nopagebreak[4]%
    \begin{list}{}{\leftmargin=0pt}
    \item[]
}{%
    \end{list}
}

% ═══════════════════════════════════════════════════════════════
\begin{document}

% Kopfbereich
\begin{center}
    {\Large\bfseries\textcolor{\fachfarbe}{Arbeitsblatt: \thema}}
\end{center}

\vspace{5pt}

\noindent
\begin{tabularx}{\textwidth}{|X|X|X|}
\hline
\textbf{Name:} \luecke{4cm} &
\textbf{Klasse:} \klasse &
\textbf{Datum:} \luecke{2cm} \\
\hline
\end{tabularx}

\vspace{15pt}

% ═══════════════════════════════════════════════════════════════
% MERKKASTEN 1: Verb querer ($\rightarrow$ FZ1)
% ═══════════════════════════════════════════════════════════════

\begin{merkkasten}[Merkkasten: Das Verb \textit{querer} (wollen) -- Stammwechsel e $\rightarrow$ ie]

\begin{center}
\begin{tabular}{|l|l|l|l|}
\hline
\textbf{Person} & \textbf{Pronomen} & \textbf{querer} & \textbf{Stammwechsel} \\
\hline
1. Sg. & yo & \luecke{2cm} & e $\rightarrow$ ie \\
\hline
2. Sg. & tú & \luecke{2cm} & e $\rightarrow$ ie \\
\hline
3. Sg. & él/ella & \luecke{2cm} & e $\rightarrow$ ie \\
\hline
1. Pl. & nosotros/-as & \luecke{2cm} & -- \\
\hline
\end{tabular}
\end{center}

\vspace{0.5em}
\textit{Übertrage die Formen aus dem Lernpfad oder von der Tafel!}
\end{merkkasten}

\vspace{10pt}

% ═══════════════════════════════════════════════════════════════
% MERKKASTEN 2: Verb poder ($\rightarrow$ FZ1)
% ═══════════════════════════════════════════════════════════════

\begin{merkkasten}[Merkkasten: Das Verb \textit{poder} (können) -- Stammwechsel o $\rightarrow$ ue]

\begin{center}
\begin{tabular}{|l|l|l|l|}
\hline
\textbf{Person} & \textbf{Pronomen} & \textbf{poder} & \textbf{Stammwechsel} \\
\hline
1. Sg. & yo & \luecke{2cm} & o $\rightarrow$ ue \\
\hline
2. Sg. & tú & \luecke{2cm} & o $\rightarrow$ ue \\
\hline
3. Sg. & él/ella & \luecke{2cm} & o $\rightarrow$ ue \\
\hline
1. Pl. & nosotros/-as & \luecke{2cm} & -- \\
\hline
\end{tabular}
\end{center}

\vspace{0.5em}
\textit{Beachte: Bei nosotros/vosotros gibt es KEINEN Stammwechsel!}
\end{merkkasten}

\vspace{15pt}

% ═══════════════════════════════════════════════════════════════
% AUFGABE 1: Konjugation zuordnen ($\rightarrow$ FZ1, AFB I)
% ═══════════════════════════════════════════════════════════════

\begin{aufgabe}{1}
Verbinde die Pronomen mit der richtigen Form. \punktebox{6}
\end{aufgabe}

\begin{center}
\begin{tabular}{rcccl}
\textbf{querer:} & & & & \\[0.5em]
yo & \luecke{3cm} & & \luecke{3cm} & quiere \\[0.8em]
tú & \luecke{3cm} & & \luecke{3cm} & quiero \\[0.8em]
él/ella & \luecke{3cm} & & \luecke{3cm} & quieres \\[1em]
\textbf{poder:} & & & & \\[0.5em]
yo & \luecke{3cm} & & \luecke{3cm} & puede \\[0.8em]
tú & \luecke{3cm} & & \luecke{3cm} & puedo \\[0.8em]
él/ella & \luecke{3cm} & & \luecke{3cm} & puedes \\
\end{tabular}
\end{center}

\vspace{15pt}

% ═══════════════════════════════════════════════════════════════
% AUFGABE 2: Lückentext ($\rightarrow$ FZ3, AFB II)
% ═══════════════════════════════════════════════════════════════

\begin{aufgabe}{2}
Setze die richtige Form von \textit{querer} oder \textit{poder} ein. \punktebox{6}
\end{aufgabe}

\begin{hinweis}
\textbf{Wortbank:} quiero | quieres | quiere | queremos | puedo | puedes | puede | podemos
\end{hinweis}

\vspace{0.5em}

\begin{enumerate}[label=\alph*)]
    \item Yo \luecke{2.5cm} comprar una camiseta roja. (wollen)
    \item ¿\luecke{2.5cm} ver esa falda azul? (können, ich)
    \item Él \luecke{2.5cm} unos pantalones negros. (wollen)
    \item Nosotros \luecke{2.5cm} pagar con tarjeta. (können)
    \item ¿Tú \luecke{2.5cm} probarte los zapatos? (wollen)
    \item María no \luecke{2.5cm} encontrar su talla. (können)
\end{enumerate}

\vspace{15pt}

% ═══════════════════════════════════════════════════════════════
% AUFGABE 3: Dialog vervollständigen ($\rightarrow$ FZ2, AFB II)
% ═══════════════════════════════════════════════════════════════

\begin{aufgabe}{3}
Vervollständige den Einkaufsdialog. Nutze die Redemittel aus dem Kasten. \punktebox{6}
\end{aufgabe}

\begin{hinweis}
\textbf{Redemittel:} Quiero... | ¿Puedo ver...? | ¿Puedo probármelo/la? | ¿Cuánto cuesta? | ¿Puedo pagar con tarjeta?
\end{hinweis}

\vspace{0.5em}

\noindent
\textbf{Vendedor:} \es{¡Buenos días! ¿En qué puedo ayudarle?}

\noindent
\textbf{Cliente:} \luecke{8cm} (ich will ein T-Shirt)

\noindent
\textbf{Vendedor:} \es{¿De qué color?}

\noindent
\textbf{Cliente:} \es{Rojo, por favor.} \luecke{5cm} (kann ich es anprobieren?)

\noindent
\textbf{Vendedor:} \es{¡Claro! El probador está allí.}

\noindent
\textbf{Cliente:} \es{Me queda bien.} \luecke{4cm} (wie viel kostet es?)

\noindent
\textbf{Vendedor:} \es{Son 25 euros.}

\noindent
\textbf{Cliente:} \luecke{6cm} (kann ich mit Karte bezahlen?)

\noindent
\textbf{Vendedor:} \es{Sí, por supuesto.}

\vspace{15pt}

% ═══════════════════════════════════════════════════════════════
% AUFGABE 4: Freier Dialog ($\rightarrow$ FZ4, AFB III)
% ═══════════════════════════════════════════════════════════════

\begin{aufgabe}{4}
Schreibe einen eigenen Einkaufsdialog (mind. 8 Zeilen). Verwende \textit{querer} und \textit{poder}! \punktebox{6}
\end{aufgabe}

\begin{hinweis}
Verwende: Begrüßung, Wunsch äußern, nach Preis fragen, bezahlen. Bonus: Erkläre die Regel für den Stammwechsel!
\end{hinweis}

\noindent\rule{\textwidth}{0.4pt}\vspace{8pt}
\noindent\rule{\textwidth}{0.4pt}\vspace{8pt}
\noindent\rule{\textwidth}{0.4pt}\vspace{8pt}
\noindent\rule{\textwidth}{0.4pt}\vspace{8pt}
\noindent\rule{\textwidth}{0.4pt}\vspace{8pt}
\noindent\rule{\textwidth}{0.4pt}\vspace{8pt}
\noindent\rule{\textwidth}{0.4pt}\vspace{8pt}
\noindent\rule{\textwidth}{0.4pt}\vspace{8pt}

\vspace{10pt}

% ═══════════════════════════════════════════════════════════════
% BONUSBOX: Stammwechsel-Regel ($\rightarrow$ FZ4, Erweiterung)
% ═══════════════════════════════════════════════════════════════

\begin{merkkasten}[Bonus: Die Stammwechsel-Regel]
\textbf{Wann wechselt der Stammvokal?}

\vspace{0.5em}
\luecke{12cm}

\vspace{0.5em}
\textbf{Weitere Verben mit Stammwechsel:}

e $\rightarrow$ ie: \luecke{6cm}

o $\rightarrow$ ue: \luecke{6cm}
\end{merkkasten}

\vspace{20pt}
\hrule
\vspace{10pt}

\begin{flushright}
\textbf{Punkte:} \luecke{1cm} / \maxpunkte
\end{flushright}

\end{document}
